\section{Tinjauan Pustaka}

\subsection{Gramian Angular Field}
GAF mengubah time series $X=(x_1,\dots,x_n)$ menjadi matriks $n \times n$~\cite{wang2015}. Data terlebih dahulu diskalakan ke rentang $[-1,1]$, lalu tiap nilai diubah menjadi sudut $\phi_i=\arccos(\tilde{x}_i)$. Gramian Angular Summation Field (GASF) dan Gramian Angular Difference Field (GADF) didefinisikan sebagai
\begin{equation}
\mathrm{GASF}_{ij}=\cos(\phi_i+\phi_j), \qquad \mathrm{GADF}_{ij}=\sin(\phi_i-\phi_j).
\end{equation}
Setiap piksel $(i,j)$ menggambarkan hubungan antara titik waktu ke-$i$ dan ke-$j$, sehingga satu titik yang berubah akan mengubah satu baris dan satu kolom pada gambar.

\subsection{Markov Transition Field}
MTF membagi nilai time series ke dalam $Q$ kelompok (bin), menghitung peluang perpindahan dari satu bin ke bin lain pada titik waktu berurutan, lalu menyusun matriks $n\times n$ dengan $M_{ij}=w_{q(i)\,q(j)}$, yaitu peluang transisi antara bin titik ke-$i$ dan bin titik ke-$j$~\cite{wang2015}. Pustaka pyts~\cite{faouzi2020} menyediakan tiga cara pembagian bin: \textit{uniform} (lebar bin sama), \textit{quantile} (jumlah titik per bin sama), dan \textit{normal} (batas bin mengikuti kuantil distribusi normal baku).

\subsection{CNN dan Grad-CAM}
CNN mengenali pola pada gambar melalui lapisan konvolusi, pooling, dan lapisan terhubung penuh. Grad-CAM~\cite{selvaraju2017} menjelaskan keputusan CNN dengan menghitung bobot tiap peta fitur pada lapisan konvolusi terakhir dari rata-rata gradien keluaran terhadap peta fitur tersebut, lalu menjumlahkan peta fitur dengan bobot itu dan membuang nilai negatif. Hasilnya adalah peta panas yang menunjukkan area gambar yang paling mendorong keputusan.

\subsection{Statistik Matriks dan Korelasi Pearson}
Rahadian dkk.~\cite{rahadian2023} menghitung rata-rata (mean), simpangan baku (SD), dan jumlah nilai unik (count) dari setiap matriks hasil encoding. Nilai-nilai ini dirata-ratakan untuk seluruh data normal dan seluruh data anomali, lalu selisih absolutnya disebut \dmean, \dsd, dan \dcount. Kekuatan hubungan antara tiap selisih dan akurasi CNN diukur dengan koefisien korelasi Pearson $r$, yang bernilai antara $-1$ dan $1$.
